\documentclass{article}
\usepackage{mathpazo}
\usepackage{graphicx}
\usepackage{subfigure}
\usepackage[latin1]{inputenc}
\usepackage{geometry}

\usepackage{fancyhdr}
\newcommand*\proptitle{Leveraging Everyday Usage of Programs to Eliminate Bugs}
\pagestyle{fancy}
\lhead{\proptitle}
\rhead{}
\renewcommand{\footrulewidth}{0.4pt}
\lfoot{Baris Kasikci (University of Michigan)}
\cfoot{}
\rfoot{\thepage}
% The NSF grant proposal guide is clear, margins must be 1 in on all sides.
\geometry{verbose, papersize={8.5in,11in}, margin=1in, headheight=14pt}

\clubpenalty10000
\widowpenalty10000
\hyphenpenalty5000

\begin{document}

%\renewcommand*\thepage{I-\arabic{page}}
\begin{center}
\textbf{\Large Facilities, Equipment, and Other Resources -- University of Michigan}
\end{center}

\vspace{1em}

\subsection*{Computer Facilities}

The facilities used to support the proposed work at the University of Michigan are located primarily in the Software Systems Lab (SSL) and the Computer Engineering Lab (CE Lab). CE lab provides state-of-the-art simulation facilities thanks equipment grants from AMD, Intel Corporation, and Sun Microsystems. The SSL and CE Lab have access to a 200+ machine collective server pool, which we will use in the context of this project to develop the distributed server system that aggregates user information and analyzes it.
 
SSL has been established by groups of faculty members of the Computer Science and Engineering Division of the EECS department at Michigan who share a common interest in experimental software systems. SSL houses several research testbeds that are used to support research in operating systems, security, fault-tolerant systems, networked and distributed systems, and databases. The testbeds contain equipment for conducting, controlling, and monitoring experiments, including gigabit networks, network-accessible power switches, uninterruptible power supplies, and remote KVM switches. The EECS department provides machine-room space for the testbeds and office space for graduate student researchers.

This project will also have access to a high-performance computing cluster called FLUX managed by the College of Engineering for large-scale application tests.  Flux is provided by the Center for Advanced Research Computing at the University of Michigan and operated by the College of Engineering's High Performance Computing Group within the Computer Aided Engineering Network (CAEN). The cluster is housed in the Modular Data Center (MDC) near the CSE Building.  FLUX staff have over 15 years of operational experience supporting high-performance computing environments. FLUX currently consists of roughly 16,000 cores; 10,000 FLUX cores can be accessed by all UM researchers via a subscription model, while the remaining 6,000 cores support 16 distinct research groups through externally funded research. In the past year, FLUX's five systems support staff and five user support staff have supported the delivery of over 16,000 CPU-years of computing time and over 7 million compute jobs. This project will use FLUX to simulate an in-production environment and in particular user executions (referred to as clients in the project description).

The PI has had access to upcoming hardware technologies (e.g., Intel Optane DC Persistent Memory, Intel HARP, etc.) through Intel vLab, which is an online portal for University collaboration. In the context of this proposal, the PI will continue his access to vLab in collaboration with Intel to explore upcoming features in Intel Processor Trace (see attached letter from Dr. Pokam at Intel).

Regarding the educational plan, the PI has already obtained a test board for Microbit from Microsoft and will seek accesses to future boards (see attached letter from Dr. Tom Ball at Microsoft).

\subsection*{Undergraduate Projects Lab}

The undergraduate projects laboratory is a facility in the CSE department where undergraduate students can work on projects in the hardware/software area.  This space will house the majority of the undergraduates that will join this project.

\subsection*{Office Space}

The PI and GSRAs have access to adequate office space within the CSE Building in the EECS Department.  This space includes conference rooms for collaborative activities and machine rooms for locating computer systems as needed.  

\subsection*{Administrative Support}

University of Michigan will make available appropriate administrative, financial and
computing support staff resources to assist the PI and graduate students
involved in this research endeavor.

\end{document}



